% T-tipi faz kolu için ortak vektör çizim tabanı.
% Argümanlar: Q1, Q2, Q3, Q4 çizgi renkleri ve ek TikZ yolu.
\newcommand{\ttypevector}[5]{%
\begin{tikzpicture}[
  x=0.76cm,y=0.56cm,
  terminal/.style={circle,fill=black,inner sep=1.25pt},
  sw/.style={rounded corners=1pt,minimum width=7mm,minimum height=8mm,fill=white,font=\scriptsize},
  bus/.style={line width=0.8pt},
  current/.style={-{Latex[length=2mm]},red!80!black,line width=1.7pt},
  activepath/.style={blue!70!black,line width=2.0pt}
]
  % Bölünmüş DA bara ve kapasitörler
  \draw[bus] (0,5) -- (4,5);
  \draw[bus] (0,-1) -- (4,-1);
  \draw[bus] (0,5) -- (0,3.65);
  \draw[bus] (-0.32,3.65) -- (0.32,3.65);
  \draw[bus] (-0.32,3.35) -- (0.32,3.35);
  \draw[bus] (0,3.35) -- (0,2);
  \draw[bus] (0,2) -- (0,0.65);
  \draw[bus] (-0.32,0.65) -- (0.32,0.65);
  \draw[bus] (-0.32,0.35) -- (0.32,0.35);
  \draw[bus] (0,0.35) -- (0,-1);
  \node[font=\scriptsize,left] at (0,5) {$+V_{\mathrm{dc}}/2$};
  \node[font=\scriptsize,left] at (0,2) {$N$};
  \node[font=\scriptsize,left] at (0,-1) {$-V_{\mathrm{dc}}/2$};
  \node[font=\scriptsize,left] at (-0.38,3.5) {$C_A$};
  \node[font=\scriptsize,left] at (-0.38,0.5) {$C_B$};

  % Ana anahtarlar ve çift yönlü nötr kol
  \node[sw,draw=#1,line width=1.25pt] (q1) at (4,3.65) {$Q_1$};
  \node[sw,draw=#2,line width=1.25pt] (q2) at (4,0.35) {$Q_2$};
  \node[sw,draw=#3,line width=1.25pt] (q3) at (1.55,2) {$Q_3$};
  \node[sw,draw=#4,line width=1.25pt] (q4) at (2.75,2) {$Q_4$};
  \draw[bus] (4,5) -- (q1.north);
  \draw[bus] (q1.south) -- (4,2);
  \draw[bus] (4,2) -- (q2.north);
  \draw[bus] (q2.south) -- (4,-1);
  \draw[bus] (0,2) -- (q3.west);
  \draw[bus] (q3.east) -- (q4.west);
  \draw[bus] (q4.east) -- (5.15,2);
  \node[terminal] at (0,2) {};
  \node[terminal] (out) at (4,2) {};
  \draw[-{Latex[length=1.8mm]},thick] (out) -- (5.15,2);
  \node[font=\scriptsize,right] at (5.15,2) {$v_{xN},\ i_x$};
  #5
\end{tikzpicture}%
}

% Faz başına tam T-tipi güç devresi: bölünmüş DA bara, dört IGBT ve LCL filtre.
% Eski tezde C olarak gösterilen filtre kapasitörü, modelle tutarlılık için C_f olarak adlandırılmıştır.
\newcommand{\ttypefulltopology}{%
\begin{tikzpicture}[
  x=0.88cm,y=0.68cm,
  wire/.style={line width=0.9pt},
  coilpath/.style={decorate,decoration={coil,aspect=0.48,segment length=2.8mm,amplitude=1.8mm}},
  igbt/.style={circle,draw,line width=0.9pt,minimum size=10mm,inner sep=0pt,fill=white},
  terminal/.style={circle,fill=black,inner sep=1.35pt},
  label/.style={font=\scriptsize,align=center}
]
  % Bölünmüş DA bara ve orta nokta
  \draw[wire] (0,5) -- (5.3,5);
  \draw[wire] (0,-1) -- (5.3,-1);
  \draw[wire] (0,5) -- (0,3.65);
  \draw[wire] (-0.38,3.65) -- (0.38,3.65);
  \draw[wire] (-0.38,3.35) -- (0.38,3.35);
  \draw[wire] (0,3.35) -- (0,2);
  \draw[wire] (0,2) -- (0,0.65);
  \draw[wire] (-0.38,0.65) -- (0.38,0.65);
  \draw[wire] (-0.38,0.35) -- (0.38,0.35);
  \draw[wire] (0,0.35) -- (0,-1);
  \node[label,left] at (0,5) {$\mathrm{DC}^{+}$};
  \node[label,left] at (0,2) {$N$};
  \node[label,left] at (0,-1) {$\mathrm{DC}^{-}$};
  \node[label,left=1mm] at (-0.38,3.50) {$C_A$};
  \node[label,left=1mm] at (-0.38,0.50) {$C_B$};
  \node[terminal] at (0,2) {};

  % Ana faz kolu IGBT'leri. D_1 ve D_2 antiparalel diyotları belirtir.
  \node[igbt] (q1) at (5.3,3.65) {};
  \draw[wire] ($(q1.center)+(0,-0.30)$) -- ($(q1.center)+(0,0.30)$);
  \draw[wire] ($(q1.center)+(-0.30,-0.18)$) -- ($(q1.center)+(-0.12,-0.18)$);
  \draw[wire] ($(q1.center)+(-0.30,-0.18)$) -- ($(q1.center)+(-0.30,0.18)$);
  \node[label,right=1.5mm] at (q1.east) {$Q_1$\\IGBT+$D_1$};
  \node[igbt] (q2) at (5.3,0.35) {};
  \draw[wire] ($(q2.center)+(0,-0.30)$) -- ($(q2.center)+(0,0.30)$);
  \draw[wire] ($(q2.center)+(-0.30,-0.18)$) -- ($(q2.center)+(-0.12,-0.18)$);
  \draw[wire] ($(q2.center)+(-0.30,-0.18)$) -- ($(q2.center)+(-0.30,0.18)$);
  \node[label,right=1.5mm] at (q2.east) {$Q_2$\\IGBT+$D_2$};
  \draw[wire] (5.3,5) -- (q1.north);
  \draw[wire] (q1.south) -- (5.3,2);
  \draw[wire] (5.3,2) -- (q2.north);
  \draw[wire] (q2.south) -- (5.3,-1);
  \node[terminal] (pole) at (5.3,2) {};

  % Orta noktayı faz kutbuna bağlayan ters seri IGBT çifti
  \node[igbt] (q3) at (2.05,2) {};
  \draw[wire] ($(q3.center)+(-0.30,0)$) -- ($(q3.center)+(0.30,0)$);
  \draw[wire] ($(q3.center)+(-0.18,-0.30)$) -- ($(q3.center)+(-0.18,-0.12)$);
  \node[label,above=1.2mm] at (q3.north) {$Q_3$ (IGBT+$D_3$)};
  \node[igbt] (q4) at (3.55,2) {};
  \draw[wire] ($(q4.center)+(-0.30,0)$) -- ($(q4.center)+(0.30,0)$);
  \draw[wire] ($(q4.center)+(0.18,-0.30)$) -- ($(q4.center)+(0.18,-0.12)$);
  \node[label,below=1.2mm] at (q4.south) {$Q_4$ (IGBT+$D_4$)};
  \draw[wire] (0,2) -- (q3.west);
  \draw[wire] (q3.east) -- (q4.west);
  \draw[wire] (q4.east) -- (pole);

  % Faz başına LCL filtre ve şebeke/yük terminali
  \coordinate (l1a) at (6.35,2);
  \coordinate (l1b) at (7.45,2);
  \coordinate (cnode) at (8.15,2);
  \coordinate (l2a) at (8.85,2);
  \coordinate (l2b) at (9.95,2);
  \draw[wire] (pole) -- (l1a);
  \draw[coilpath,line width=0.9pt] (l1a) -- (l1b);
  \draw[wire] (l1b) -- (cnode) -- (l2a);
  \draw[coilpath,line width=0.9pt] (l2a) -- (l2b);
  \draw[wire,-{Latex[length=1.8mm]}] (l2b) -- (11.1,2);
  \node[label,above=1.5mm] at ($(l1a)!0.5!(l1b)$) {$L_1$};
  \node[label,above=1.5mm] at ($(l2a)!0.5!(l2b)$) {$L_2$};
  \node[label,right] at (11.1,2) {$v_g,\ i_2$};
  \node[terminal] at (cnode) {};
  \draw[wire] (cnode) -- (8.15,1.25);
  \draw[wire] (7.78,1.25) -- (8.52,1.25);
  \draw[wire] (7.78,0.95) -- (8.52,0.95);
  \draw[wire] (8.15,0.95) -- (8.15,-1.45) -- (-0.85,-1.45);
  \draw[wire] (-0.85,-1.45) -- (-0.85,2) -- (0,2);
  \node[label,right=1mm] at (8.52,1.10) {$C_f$ (eski $C$)};
  \draw[-{Latex[length=1.8mm]},red!80!black,line width=1.1pt] (6.0,2.30) -- (7.0,2.30);
  \node[label,above] at (6.5,2.30) {$i_1$};
\end{tikzpicture}%
}

% Senkron-buck eşdeğeri: üst ray, alt ray, üst ve alt anahtar etiketleri.
\newcommand{\syncbuckvector}[4]{%
\begin{tikzpicture}[
  x=0.82cm,y=0.65cm,
  sw/.style={draw,rounded corners=1pt,minimum width=8mm,minimum height=8mm,fill=white,font=\scriptsize},
  coilpath/.style={decorate,decoration={coil,aspect=0.48,segment length=2.7mm,amplitude=1.8mm}},
  current/.style={-{Latex[length=2mm]},red!80!black,line width=1.4pt}
]
  \draw[line width=0.8pt] (0,3.2) -- (1.8,3.2);
  \draw[line width=0.8pt] (0,0) -- (1.8,0);
  \node[font=\scriptsize,left] at (0,3.2) {#1};
  \node[font=\scriptsize,left] at (0,0) {#2};
  \node[sw,draw=green!55!black,line width=1.2pt] (qh) at (1.8,2.45) {#3};
  \node[sw,draw=gray!65,line width=1.2pt] (ql) at (1.8,0.75) {#4};
  \draw[line width=0.8pt] (1.8,3.2) -- (qh.north);
  \draw[line width=0.8pt] (qh.south) -- (1.8,1.6);
  \draw[line width=0.8pt] (1.8,1.6) -- (ql.north);
  \draw[line width=0.8pt] (ql.south) -- (1.8,0);
  \fill (1.8,1.6) circle (1.5pt);
  \coordinate (l1a) at (2.45,1.6);
  \coordinate (l1b) at (3.35,1.6);
  \coordinate (cnode) at (3.88,1.6);
  \coordinate (l2a) at (4.30,1.6);
  \coordinate (l2b) at (5.20,1.6);
  \draw[line width=0.8pt] (1.8,1.6) -- (l1a);
  \draw[coilpath,line width=0.9pt] (l1a) -- (l1b);
  \draw[line width=0.8pt] (l1b) -- (cnode) -- (l2a);
  \draw[coilpath,line width=0.9pt] (l2a) -- (l2b);
  \draw[current] (l2b) -- (6.0,1.6);
  \node[font=\scriptsize,above=2mm] at ($(l1a)!0.5!(l1b)$) {$L_1$};
  \node[font=\scriptsize,above=2mm] at ($(l2a)!0.5!(l2b)$) {$L_2$};
  \node[font=\scriptsize,right] at (6.0,1.6) {$v_g,\ i_L$};
  \fill (cnode) circle (1.3pt);
  \draw[line width=0.8pt] (cnode) -- (3.88,1.0);
  \draw[line width=0.8pt] (3.55,1.0) -- (4.21,1.0);
  \draw[line width=0.8pt] (3.55,0.75) -- (4.21,0.75);
  \draw[line width=0.8pt] (3.88,0.75) -- (3.88,0.2);
  \draw[line width=0.8pt] (3.58,0.2) -- (4.18,0.2);
  \draw[line width=0.8pt] (3.68,0.05) -- (4.08,0.05);
  \node[font=\scriptsize,right] at (4.18,0.88) {$C_f$};
  \draw[blue!70!black,line width=1.8pt] (0.5,3.2) -- (qh.north);
  \draw[blue!70!black,line width=1.8pt] (qh.south) -- (1.8,1.6) -- (l1a);
\end{tikzpicture}%
}
